\begin{tomtat}
	\begin{minipage}[t]{0.3\textwidth}
		\vspace{0pt}
		\begin{tomtat}
			\begin{center}
				\textbf{ĐỀ 26}
			\end{center}
			
		\end{tomtat}
	\end{minipage}
	\hfill
	\begin{minipage}[t]{0.5\textwidth}
		\vspace{0pt}
		\begin{tomtat}
			\begin{center}
				\includegraphics[scale=0.08]{images/26}
			\end{center}
		\end{tomtat}
	\end{minipage}
\end{tomtat}
\textbf{I. PHẦN TRẮC NGHIỆM}
\begin{ex}
	Hệ phương trình $\begin{cases} 2x+y=8 \\ x-y=1 \end{cases}$ có nghiệm $(x;y)$ là:
	\choice{$(-3; 2)$}{$(3; -2)$}{$(3; 2)$}{$( -3; -2)$}
	\loigiai{
		Ta có hệ phương trình:
		$$ \begin{cases} 2x+y=8 \quad &(1) \\ x-y=1 \quad &(2) \end{cases} $$
		Cộng vế theo vế của phương trình (1) và (2), ta được:
		$(2x+y) + (x-y) = 8+1$
		$3x = 9$
		$x = \frac{9}{3}$
		$x = 3$.
		Thay $x=3$ vào phương trình (2):
		$3-y=1$
		$y = 3-1$
		$y = 2$.
		Vậy nghiệm của hệ phương trình là $(x; y) = (3; 2)$.
		\\ \textbf{Đáp Án C}
	}
\end{ex}

\begin{ex}
	Hai số có tổng bằng 3 và tích bằng –10 là nghiệm của phương trình:
	\choice{$x^2-3x-10=0$}{$x^2-3x+10=0$}{$x^2+3x-10=0$}{$x^2+3x+10=0$}
	\loigiai{
		Nếu hai số có tổng là S và tích là P thì chúng là nghiệm của phương trình $x^2 - Sx + P = 0$.
		Theo đề bài, $S = 3$ và $P = -10$.
		Vậy phương trình cần tìm là: $x^2 - 3x + (-10) = 0 \Leftrightarrow x^2-3x-10=0$.
		\\ \textbf{Đáp Án A}
	}
\end{ex}

\begin{ex}
	Xác định đường thẳng $y=ax+b \; (a \neq 0)$ có hệ số gốc bằng -4 và đi qua điểm A(3;-2).
	\choice{$y=-4x+10$}{$y=4x+10$}{$y=-4x-10$}{$y=-4x$}
	\loigiai{
		Đường thẳng có hệ số góc bằng -4 nên $a=-4$.
		Phương trình đường thẳng có dạng $y = -4x+b$.
		Đường thẳng đi qua điểm A(3; -2), nên tọa độ của A thỏa mãn phương trình đường thẳng:
		$-2 = -4(3) + b$
		$-2 = -12 + b$
		$b = -2 + 12$
		$b = 10$.
		Vậy phương trình đường thẳng là $y=-4x+10$.
		\\ \textbf{Đáp Án A}
	}
\end{ex}

\begin{ex}
	Điểm $P(-1;-2)$ thuộc đồ thị hàm số $y=mx^2$ khi m bằng:
	\choice{$2$}{$-2$}{$4$}{$-4$}
	\loigiai{
		Điểm $P(-1;-2)$ thuộc đồ thị hàm số $y=mx^2$, nên tọa độ của P thỏa mãn phương trình hàm số:
		$-2 = m(-1)^2$
		$-2 = m \cdot 1$
		$m = -2$.
		\\ \textbf{Đáp Án B}
	}
\end{ex}

\begin{ex}
	Cho tam giác ABC vuông tại A có $AB=3$cm, $AC=4$cm. Khi đó $\tan C$ bằng:
	\choice{$\frac{3}{4}$}{$\frac{3}{5}$}{$\frac{4}{3}$}{$\frac{4}{5}$}
	\loigiai{
		Trong tam giác ABC vuông tại A:
		$\tan C = \frac{\text{cạnh đối}}{\text{cạnh kề}} = \frac{AB}{AC}$.
		Với $AB=3$cm và $AC=4$cm:
		$\tan C = \frac{3}{4}$.
		\\ \textbf{Đáp Án A}
	}
\end{ex}
\begin{ex}
	Ngọn hải đăng Tiên Nữ cao 22,1m được xây dựng năm 2000 tại đảo Tiên Nữ thuộc quần đảo Trường Sa, huyện Trường Sa, tỉnh Khánh Hòa. Ngoài nhiệm vụ đảm bảo an toàn hàng hải trong khu vực quần đảo, ngọn hải đăng này còn là cột mốc chủ quyền của Tổ quốc trên Biển Đông. Một con tàu nhìn thấy ngọn hải đăng Tiên Nữ theo một góc là $\alpha = 1^\circ15'$. Hỏi tàu cách ngọn hải đăng bao nhiêu mét? (Làm tròn kết quả đến chữ số thập phân thứ nhất).
	\choice{$1012,8$}{$1013,1$}{$1012,7$}{$1012,9$}
	\loigiai{
		Gọi chiều cao của ngọn hải đăng là $h = 22,1$ m.
		Gọi khoảng cách từ tàu đến chân ngọn hải đăng là $d$ (mét).
		Góc nhìn từ tàu đến đỉnh ngọn hải đăng là $\alpha = 1^\circ15'$.
		Đổi $1^\circ15' = 1 + \frac{15}{60} = 1 + \frac{1}{4} = 1,25^\circ$.
		Ta có mối quan hệ lượng giác trong tam giác vuông tạo bởi tàu, chân hải đăng và đỉnh hải đăng:
		$\tan(\alpha) = \frac{h}{d}$
		$d = \frac{h}{\tan(\alpha)} = \frac{22,1}{\tan(1,25^\circ)}$.
		Sử dụng máy tính:
		$\tan(1,25^\circ) \approx 0,0218196$
		$d \approx \frac{22,1}{0,0218196} \approx 1012,849...$
		Làm tròn kết quả đến chữ số thập phân thứ nhất, ta được $d \approx 1012,8$ mét.
		\\ \textbf{Đáp Án A}
	}
\end{ex}

\begin{ex}
	Một bồn nước inox có dạng hình trụ chiều cao 2m, bán kính đáy 0,8m. Hỏi bồn nước này đựng được bao nhiêu lít nước (Bỏ qua bề dày thành bể, làm tròn kết quả đến chữ số thập phân thứ nhất và lấy $\pi \approx 3,14$).
	\choice{$4001$}{$4000$}{$101$}{$4019$}
	\loigiai{
		Chiều cao của hình trụ $h = 2$ m.
		Bán kính đáy của hình trụ $r = 0,8$ m.
		Thể tích của hình trụ được tính bằng công thức $V = \pi r^2 h$.
		$V = \pi \cdot (0,8)^2 \cdot 2$
		$V = \pi \cdot 0,64 \cdot 2$
		$V = 1,28\pi \text{ (m}^3\text{)}$.
		Với $\pi \approx 3,14$:
		$V \approx 1,28 \cdot 3,14 = 4,0192 \text{ (m}^3\text{)}$.
		Đổi sang lít: $1 \text{ m}^3 = 1000$ lít.
		$V \approx 4,0192 \cdot 1000 = 4019,2$ lít.
		Đề bài yêu cầu làm tròn kết quả đến chữ số thập phân thứ nhất, vậy kết quả là $4019,2$ lít.
		\\ \textbf{Đáp Án D}
	}
\end{ex}

\begin{ex}
	Một hộp có 30 thẻ cùng loại, mỗi thẻ được ghi một trong các số 1; 2; 3; 4; 5; ...; 29; 30; hai thẻ khác nhau thì ghi số khác nhau. Rút ngẫu nhiên một thẻ trong hộp. Xác suất của biến cố “Số xuất hiện trên thẻ được rút ra là số có hai chữ số và tổng các chữ số bằng 6” là:
	\choice{$\frac{1}{30}$}{$\frac{1}{10}$}{$\frac{1}{15}$}{$\frac{1}{2}$}
	\loigiai{
		Số phần tử của không gian mẫu là $n(\Omega) = 30$ (vì có 30 thẻ).
		Gọi A là biến cố “Số xuất hiện trên thẻ được rút ra là số có hai chữ số và tổng các chữ số bằng 6”.
		Các số có hai chữ số mà tổng các chữ số bằng 6 là:
		\begin{itemize}
			\item 15 (1+5=6)
			\item 24 (2+4=6)
			\item 33 (3+3=6) - Không có thẻ này vì các số từ 1 đến 30.
			\item 42 (4+2=6) - Không có thẻ này.
			\item 51 (5+1=6) - Không có thẻ này.
			\item 60 (6+0=6) - Không có thẻ này.
		\end{itemize}
		Trong các số từ 1 đến 30, các số có hai chữ số và tổng các chữ số bằng 6 là:
		15 (vì $10 \le 15 \le 30$ và $1+5=6$)
		24 (vì $10 \le 24 \le 30$ và $2+4=6$)
		Các số khác có tổng chữ số bằng 6 nhưng lớn hơn 30 (ví dụ: 33, 42, 51, 60) hoặc không phải số có 2 chữ số (ví dụ: 06) không được xét.
		Vậy, các kết quả thuận lợi cho biến cố A là $A = \{15; 24\}$.
		Số phần tử của biến cố A là $n(A) = 2$.
		Xác suất của biến cố A là:
		$P(A) = \frac{n(A)}{n(\Omega)} = \frac{2}{30} = \frac{1}{15}$.
		\\ \textbf{Đáp Án C}
	}
\end{ex}
\newpage
\textbf{II. PHẦN TỰ LUẬN}
\begin{ex}
	Rút gọn biểu thức $A = \left( \frac{1}{x+2\sqrt{x}} - \frac{1}{\sqrt{x}+2} \right) : \frac{1-\sqrt{x}}{x+4\sqrt{x}+4}$ với $x>0; x \neq 1$.
	\loigiai{
		Điều kiện: $x>0; x \neq 1$.
		Ta có:
		\begin{align*} A &= \left( \frac{1}{\sqrt{x}(\sqrt{x}+2)} - \frac{\sqrt{x}}{\sqrt{x}(\sqrt{x}+2)} \right) : \frac{1-\sqrt{x}}{(\sqrt{x}+2)^2} \\ &= \frac{1-\sqrt{x}}{\sqrt{x}(\sqrt{x}+2)} \cdot \frac{(\sqrt{x}+2)^2}{1-\sqrt{x}} \end{align*}
		Vì $x>0$ và $x \neq 1$, nên $1-\sqrt{x} \neq 0$ và $\sqrt{x}+2 \neq 0$.
		\begin{align*} A &= \frac{\cancel{(1-\sqrt{x})}}{\sqrt{x}\cancel{(\sqrt{x}+2)}} \cdot \frac{(\sqrt{x}+2)^{\cancel{2}}}{\cancel{(1-\sqrt{x})}} \\ &= \frac{\sqrt{x}+2}{\sqrt{x}} \end{align*}
		Vậy $A = \frac{\sqrt{x}+2}{\sqrt{x}}$.
	}
\end{ex}

\begin{ex}
	Giải hệ phương trình $\begin{cases} 2x-y=1 \\ 3x+y=4 \end{cases}$.
	\loigiai{
		Ta có hệ phương trình:
		$$ \begin{cases} 2x-y=1 \quad &(1) \\ 3x+y=4 \quad &(2) \end{cases} $$
		Cộng vế theo vế của phương trình (1) và (2), ta được:
		\begin{align*} (2x-y) + (3x+y) &= 1+4 \\ 5x &= 5 \\ x &= 1 \end{align*}
		Thay $x=1$ vào phương trình (2):
		\begin{align*} 3(1)+y &= 4 \\ 3+y &= 4 \\ y &= 4-3 \\ y &= 1 \end{align*}
		Vậy nghiệm của hệ phương trình là $(x;y) = (1;1)$.
	}
\end{ex}

\begin{ex}
	Hộp thứ nhất đựng 1 thẻ xanh, 1 thẻ đỏ và 1 thẻ vàng. Hộp thứ hai đựng 1 thẻ xanh, 1 thẻ đỏ. Các tấm thẻ có kích thước có khối lượng như nhau. Lần lượt lấy ra ngẫu nhiên từ mỗi hộp một tấm thẻ. Tính xác suất của biến cố “Trong 2 thẻ lấy ra có ít nhất 1 thẻ màu xanh”.
	\loigiai{
		Ký hiệu: Xanh (X), Đỏ (Đ), Vàng (V).
		Hộp 1: $\{X_1; Đ_1; V_1\}$
		Hộp 2: $\{X_2; Đ_2\}$
		
		Lấy ngẫu nhiên mỗi hộp 1 thẻ. Các kết quả có thể xảy ra (không gian mẫu $\Omega$):
		$\Omega = \{ (X_1;X_2); (X_1;Đ_2); (Đ_1;X_2); (Đ_1;Đ_2); (V_1;X_2); (V_1;Đ_2) \}$.
		Số phần tử của không gian mẫu là $n(\Omega) = 3 \times 2 = 6$.
		
		Gọi A là biến cố: “Trong 2 thẻ lấy ra có ít nhất 1 thẻ màu xanh”.
		Các kết quả thuận lợi cho biến cố A:
		\begin{itemize}
			\item Thẻ từ hộp 1 là xanh: $(X_1;X_2)$; $(X_1;Đ_2)$.
			\item Thẻ từ hộp 2 là xanh (và thẻ từ hộp 1 không phải xanh để không trùng lặp): $(Đ_1;X_2)$; $(V_1;X_2)$.
		\end{itemize}
		Vậy, $A = \{ (X_1;X_2); (X_1;Đ_2); (Đ_1;X_2); (V_1;X_2) \}$.
		Số phần tử của biến cố A là $n(A) = 4$.
		
		Xác suất của biến cố A là:
		$P(A) = \frac{n(A)}{n(\Omega)} = \frac{4}{6} = \frac{2}{3}$.
	}
\end{ex}

\begin{ex}
	Gọi $x_1, x_2$ là nghiệm của phương trình $-2x^2-6x-1=0$. Không giải phương trình tính giá trị của biểu thức $N=\frac{1}{x_1+3}+\frac{1}{x_2+3}$.
	\loigiai{
		Phương trình $-2x^2-6x-1=0$ có $a=-2, b=-6, c=-1$.
		$\Delta = b^2-4ac = (-6)^2 - 4(-2)(-1) = 36 - 8 = 28 > 0$.
		Phương trình có hai nghiệm phân biệt $x_1, x_2$.
		Theo hệ thức Vi-ét, ta có:
		$$ \begin{cases} x_1 + x_2 = -\frac{b}{a} = -\frac{-6}{-2} = -3 \\ x_1 x_2 = \frac{c}{a} = \frac{-1}{-2} = \frac{1}{2} \end{cases} $$
		Ta biến đổi biểu thức N:
		\begin{align*} N &= \frac{1}{x_1+3} + \frac{1}{x_2+3} \\ &= \frac{(x_2+3) + (x_1+3)}{(x_1+3)(x_2+3)} \\ &= \frac{x_1+x_2+6}{x_1x_2 + 3x_1 + 3x_2 + 9} \\ &= \frac{(x_1+x_2)+6}{x_1x_2 + 3(x_1+x_2) + 9} \end{align*}
		Thay giá trị từ hệ thức Vi-ét vào biểu thức N:
		\begin{align*} N &= \frac{(-3)+6}{\frac{1}{2} + 3(-3) + 9} \\ &= \frac{3}{\frac{1}{2} - 9 + 9} \\ &= \frac{3}{\frac{1}{2}} \\ &= 3 \cdot 2 \\ &= 6 \end{align*}
		Vậy $N=6$.
	}
\end{ex}
\begin{ex}
	Sau một buổi sinh hoạt ngoại khóa, nhóm bạn của Thư rủ nhau đi ăn kem ở một quán gần trường. Do quán mới khai trương nên có khuyến mại, bắt đầu từ que kem thứ 5, mỗi ly kem được giảm 1500 đồng so với ban đầu. Nhóm của Thư mua 9 ly kem với số tiền là 154500 đồng. Hỏi giá của mỗi ly kem ban đầu?
	\loigiai{
		Gọi giá của mỗi ly kem ban đầu là $x$ (đồng).
		Điều kiện: $x > 1500$ (vì giá sau khi giảm phải dương).
		
		Nhóm của Thư mua 9 ly kem.
		Số ly kem được mua với giá ban đầu là 4 ly (từ ly thứ 1 đến ly thứ 4).
		Số tiền mua 4 ly kem đầu là: $4x$ (đồng).
		
		Số ly kem được giảm giá là: $9 - 4 = 5$ ly (từ ly thứ 5 đến ly thứ 9).
		Giá của mỗi ly kem sau khi giảm là: $x - 1500$ (đồng).
		Số tiền mua 5 ly kem sau khi giảm giá là: $5(x - 1500)$ (đồng).
		
		Tổng số tiền mua 9 ly kem là 154500 đồng, nên ta có phương trình:
		\begin{align*} 4x + 5(x - 1500) &= 154500 \\ 4x + 5x - 7500 &= 154500 \\ 9x &= 154500 + 7500 \\ 9x &= 162000 \\ x &= \frac{162000}{9} \\ x &= 18000 \end{align*}
		Giá trị $x=18000$ thỏa mãn điều kiện $x > 1500$.
		
		Vậy, giá của mỗi ly kem ban đầu là 18000 đồng.
	}
\end{ex}
\begin{ex}
	Từ điểm M nằm ngoài đường tròn (O) kẻ hai tiếp tuyến MA, MB tới (O) (A, B là các tiếp điểm). Gọi H là giao điểm của AB với MO; E và F là các giao điểm của đường thẳng MO với đường tròn (O) (với ME < MF).
	\begin{enumerate}
		\item[1.] Chứng minh bốn điểm M, A, O, B cùng thuộc một đường tròn.
		\item[2.] Chứng minh MO $\perp$ AB và $HE \cdot HF = HM \cdot HO$.
		\item[3.] Kẻ đường kính BP của đường tròn (O). Đường thẳng MP cắt đường tròn (O) tại điểm N (N khác P). Gọi I là trung điểm của đoạn thẳng MH. Chứng minh ba điểm A, N, I thẳng hàng.
	\end{enumerate}
\end{ex}

\begin{ex}
	Cho các số thực dương x, y thỏa mãn $x+y \le 3$.
	Tìm giá trị nhỏ nhất của biểu thức $P = \frac{1}{5xy} + \frac{5}{x+2y+5}$.
\end{ex}

\setcounterref{ex}{0}
